\begin{proposition}[Rank divisors from the root Kodaira--Spencer morphism]\hypertarget{tag-01GN}{}\label{cdp:st:E16}
\StatementRecord{cdp:st:E16}{proved}{01GN}{\SeriesRef{cdp:eq:E.16.3}, \SeriesRef{cdp:st:E8}, \SeriesRef{cdp:st:E13}, \SeriesRef{cdp:st:E14}, \SeriesRef{cdp:st:E15}, \SeriesRef{cdp:st:E7}}

We use the logarithmic Kodaira--Spencer morphism at the specified degenerate fibres and the ordinary morphism on the finite smooth torus covers. For the original affine-twist family, including every $|p|>1$, its only rank defect is rank one at the order-three point. The logarithmic cusp has rank two. For the $(3k,4l)$ logarithmic transforms, on their full root covers the ranks at the two multiple values are respectively
\[
 \begin{cases}1&k=1,\\0&k>1,\end{cases}
 \qquad
 \begin{cases}2&l=1,\\0&l>1.\end{cases}                 \tag{E.16.1}\hypertarget{tag-01A2}{}\label{cdp:eq:E.16.1}
\]
Every other smooth point and the logarithmic cusp have rank two.

For the two-parameter family the logarithmic rank is one at each of its two rank-one nodal values, zero at the order-four root point, and two everywhere else, including the hexagonal fibre. Its determinant has simple coarse zeros at the two nodal values and order four in the root coordinate at the multiple value. For the Liu--Rollenske central family the logarithmic rank is two at both the cuspidal and hexagonal fibres and throughout the remaining smooth locus; at the order-four root point it is zero, and its determinant has order four there.

For a specified positive comb section $\sigma$ of degree parameter $m$, let $x,y$ be root coordinates of weights $3,4$, with $y=0$ at the order-three point. There is a nonzero homogeneous polynomial $P_\sigma(X,Y)$ of degree $4m$, uniquely determined up to an overall nonzero scalar by the marked period section, such that
\[
 \det\kappa_\sigma=y^2P_\sigma(x^4,y^3).             \tag{E.16.2}\hypertarget{tag-01A3}{}\label{cdp:eq:E.16.2}
\]
The rank-defect set is exactly its zero divisor. At any such point the rank is zero exactly when all four entries of the actual Kodaira--Spencer matrix vanish; otherwise it is one. Formula \eqref{cdp:eq:E.16.3} below computes those entries from $\sigma$. No unspecified choice of a comb section is assigned a fixed root set. The polynomial is nonzero at the hexagonal cusp, so that cusp always has rank two.
\end{proposition}


\begin{proof}

\begin{enumerate}[label={$\langle1\rangle$\arabic*.},leftmargin=*]
\item Adding a linearization constant or changing only the admissible affine lifts does not change the derivative of the marked period matrix. \ref{cdp:st:E8} therefore proves the first assertion. On a root cover of a higher logarithmic transform the original local good-reduction coordinate is $s=u^k$ at the first multiple value and $s=u^l$ at the second. The lifted period map is its actual holomorphic pullback, not a change of frame with a pole. The chain rule gives $k u^{k-1}\kappa(s)$ or $l u^{l-1}\kappa(s)$. \ref{cdp:st:E8} gives rank one at the original order-three centre and rank two at the original order-four centre. The covering derivative is nonzero off its centre, proving \eqref{cdp:eq:E.16.1} and the absence of further defects. The cusp coordinate is unchanged.

\item For the two-parameter family \ref{cdp:st:E13} proves, with its extension maps,
\[
 E=G=\OO\oplus\OO(-1),\qquad
 \kappa:\OO(-7)\longrightarrow E\otimes G
 \quad\text{on }\PP(1,4).
\]
Thus the four entries have weights $7,6,6,5$. Write $x$ for the weight-one coordinate vanishing at the root point and $y$ for the weight-four coordinate. Every monomial of these respective weights is divisible by $x^3,x^2,x^2,x$. Consequently the entire matrix vanishes at the root point, and its determinant is divisible by $x^4$. The determinant is a section of $\OO(12)$, whose coarse degree is three. At each of the two rank-one Mumford values, the limiting logarithmic derivative is the nonzero rank-one integral monodromy operator in the local torus frames. Its determinant therefore vanishes at least once, whereas its matrix has rank one. These two coarse zeros, together with the root zero of order at least four, account for all of the degree of $\OO(12)$. The determinant is nonzero because the logarithmic rank at the hexagonal cusp is two. All the stated orders are therefore exact, with no remaining zeros. This proves the two-parameter assertion without assuming a theta identity for its varying section.

\item For the Liu--Rollenske central member \ref{cdp:st:E14} proves the analogous map from $\OO(-3)$ on $\PP(1,4)$. Its entries have weights $3,2,2,1$; these spaces consist respectively of multiples of $x^3,x^2,x^2,x$. The determinant is a section of $\OO(4)$ and is a constant multiple of $x^4$. The constant is nonzero by the hexagonal local monodromy. Hence the only rank defect is the root point, where all four entries vanish. In particular this calculation determines the value of the extended logarithmic boundary at the cusp, even though its fibre is not normal crossing.

\item For a positive comb, \ref{cdp:st:E15} gives
\[
 \det V=\det G=\OO(24m-1),\qquad
 \kappa_\sigma:\OO(-5)\longrightarrow V\otimes G.
\]
The determinant is a section of $\det V\otimes\det G\otimes\OO(10)=\OO(48m+8)$. On $\PP(3,4)$ every weighted monomial $x^ay^b$ of this degree satisfies $b\equiv2\pmod3$ and $a\equiv0\pmod4$. Factoring $y^2$ therefore gives exactly the polynomial in \eqref{cdp:eq:E.16.2}, of coarse degree $4m$.

\item Here is the exact computation of that polynomial for any of the period sections in the construction. On a marked local root disc let the two columns of a holomorphic $4\times2$ matrix $U$ frame the Hodge subbundle of the canonical flat extension. Let $\nabla$ be its explicitly marked Gauss--Manin connection, and let $D$ be the local generator of $T_{\mathfrak B}(-p_0)$: $\partial_s$ away from the cusp and $q\partial_q$ at the cusp. Then
\[
 \kappa_\sigma=(\nabla_DU)\bmod\operatorname{im}U,
 \qquad
 \det\kappa_\sigma\ \text{is represented by}\
 \det(U_1,U_2,\nabla_DU_1,\nabla_DU_2).              \tag{E.16.3}\hypertarget{tag-01A4}{}\label{cdp:eq:E.16.3}
\]
The first expression uses the quotient frame of $G$ and the dual frame of $V=F^*$; it specifies every matrix entry and therefore also the rank-zero test. Changing the Hodge frame multiplies the determinant by the square of its determinant, exactly its transition in $\det V\otimes\det G\otimes T_{\mathfrak B}(-p_0)^{-2}$. Changing a root parameter gives the inverse square of the tangent-frame change. Thus the local calculations glue to the stated global section and determine $P_\sigma$ from the actual period data. In the cusp comparison of \ref{cdp:st:E7} and \ref{cdp:st:E15}, the logarithmic derivative has the fixed invertible integral leading matrix; its determinant is nonzero. This proves both the nonzero-polynomial assertion and the exact rank-divisor description.

\end{enumerate}
\end{proof}